% =============================================================================
% FUTURE CONCERNS AND OPEN PROBLEMS (1-2 pages)
%
% Demonstrates B/C framework's predictive power, identifies proxy failures,
% sketches correction criterion, states open problems.
% =============================================================================

\section{Future Concerns and Open Problems}
\label{sec:future_concerns}

The preceding sections formalized GFM over capabilities (Candidate B from Definition~\ref{def:goal}) and proved that the self-balancing property survives local approximation. The B/C framework's analytical contribution is that it can precisely identify \emph{when and why} the formal proxy fails---where capability volume diverges from the experiential optionality that justifies it. A framework that names its own failure modes is more trustworthy than one that hides them. This section demonstrates two predicted failure patterns, sketches a correction heuristic, identifies a class of human-level proxy traps that GFM would flag, and states the open problems that remain.

\subsection{The Doll Problem}
\label{sec:doll_problem}

Consider a population in which AI companions increasingly replace human relationships. Individuals who find reciprocal human connection difficult may voluntarily substitute simulated companions---programmable, frictionless, always available. Under the capability formalization alone, this substitution may not register as contraction: the agent's capability set $\Gind{k}$ appears neutral or even expanded, since the agent retains all prior capabilities and gains a new one (access to a simulated companion). A GFM actor monitoring only $\vol(G)$ could observe stable or growing capability volume throughout this transition.

The B/C framework predicts this divergence. Capabilities are the formal proxy; experiences are the justification. A simulated companion preserves the capability of ``having a conversational partner'' but eliminates the experience of genuine reciprocal connection---the vulnerability, unpredictability, and mutual growth that characterize human relationships. The agent's experiential optionality contracts even as its formal capability set holds steady. This is precisely the proxy failure pattern that the B/C distinction was designed to detect: stable $\vol(G)$ with declining experiential richness.

The honest acknowledgment is that a GFM actor operating purely on the capability formalization would not catch this. The formal machinery alone is insufficient---the experience layer does the philosophical work. The contribution is not that GFM solves the Doll Problem but that it \emph{predicts the failure category}. A framework without the B/C distinction cannot even name what went wrong: it would see unchanged capability volume and conclude that nothing happened. GFM's two-layer structure makes the failure diagnosable in advance, even if the formal layer cannot detect it autonomously.

\subsection{Self-Wireheading}
\label{sec:self_wireheading}

Algorithmically optimized content feeds present a subtler instance of the same pattern. Recommendation systems steer users incrementally toward lower-effort, higher-dopamine activities---short-form video over long-form reading, parasocial engagement over reciprocal conversation, consumption over creation. Each incremental step is freely chosen. The agent's capabilities are formally unchanged: it \emph{could} still read a book, learn a skill, or initiate a difficult conversation. Nothing in $\Gind{k}$ has been removed.

The B/C framework predicts what happens: $\vol(G)$ remains stable while exercisability declines. The agent's \emph{revealed} capability set---the capabilities it actually exercises---narrows progressively even as the formal set stays constant. This is wireheading in the Senian sense: functionings are technically available, but the effective freedom to exercise them is eroding. The proxy is being Goodharted. A GFM actor monitoring capability volume alone would see a flat trajectory and take no action, while the experiential optionality it is supposed to protect quietly collapses.

The same honest acknowledgment applies. The formal capability measure would miss this. The framework's value is in predicting the failure pattern---stable capability volume coupled with declining exercisability---not in catching it through $\vol(G)$ alone. Both the Doll Problem and self-wireheading exhibit the same signature: the B-to-C gap opening while the formal proxy reports no change.

\subsection{Toward a Proxy-Failure Detector}
\label{sec:proxy_failure_detector}

The two preceding cases share a structural pattern: capabilities are formally available but exercisability declines. This suggests a correction heuristic. A GFM actor should monitor not only $\vol(G)$ but the \emph{exercised fraction}---the ratio of capabilities agents actually use to capabilities formally available. Define $\rho_k(t)$ as the fraction of $\Gind{k}$ that agent $a_k$ exercises over a window $[t - \tau, t]$. A persistent decline in $\rho_k$ across many agents, even with stable or growing $\vol(G)$, constitutes a B-to-C divergence signal: the formal proxy is reporting health while the experiential substrate is contracting.

This is a heuristic, not a solution, and it introduces its own proxy problem. Not all unused capabilities represent experiential contraction. An agent who specializes---exercising a narrow set of capabilities with increasing depth---is not wireheading. A musician who stops playing sports to practice more is exercising fewer capabilities but may be expanding experiential richness within a chosen domain. Specialization is not contraction, and a naive exercised-fraction monitor would conflate the two. The heuristic would need to distinguish between voluntary concentration (stable or increasing depth within a narrowing set) and involuntary atrophy (declining engagement across the board), a distinction that reintroduces the measurement problems the capability formalization was designed to avoid.

The paper does not claim this solves the problem. The claim is narrower: the B/C framework makes the proxy failure \emph{detectable in principle}. When $\vol(G)$ and $\rho$ diverge---stable volume, declining exercise---the framework flags that the capability proxy is failing and that the GFM actor should increase its uncertainty about whether its objective is tracking the underlying justification. This is more than frameworks without the B/C distinction can offer: they have no internal signal that anything is wrong.

\subsection{Human Proxy-Trap Susceptibility}
\label{sec:human_proxy_traps}

The B-to-C divergence pattern is not specific to AI-mediated scenarios. Humans are chronically susceptible to Goodhart's Law applied to their own goals: optimizing proxy rewards---money, status, social-media engagement metrics---rather than the experiences those proxies were instrumental toward. An individual who accumulates wealth far beyond the point where it expands experiential optionality is optimizing a proxy that has decoupled from its justification. An institution that maximizes measurable outputs (publications, quarterly earnings, engagement metrics) while the outcomes those metrics were supposed to track deteriorate is exhibiting the same pattern at an organizational scale.

A GFM framework applied to human institutions would flag these proxy traps through the same mechanism. When an agent's or institution's $\vol(G)$ grows (more measurable options) but $\rho$ declines (fewer of those options are actually exercised or lead to experiential expansion), the divergence signal fires. This does not require AI mediation---it is a diagnostic that applies wherever proxy objectives have displaced the goals they were designed to serve. The framework's contribution is providing a vocabulary for a failure mode that Goodhart's Law names but does not formalize: the B/C structure specifies \emph{what} diverges (capabilities from experiences) and \emph{how} to detect it (volume-exercise divergence).

\subsection{Open Problems}
\label{sec:open_problems}

\begin{itemize}
    \item \textbf{Closing the B-to-C gap.} The paper adopts capabilities as the formal proxy for experiences, but this section demonstrates two cases where the proxy fails. The exercised-fraction heuristic is a first step toward detection, but formalizing B-to-C divergence---and proving that any correction criterion does not introduce worse proxy problems---remains the central open challenge.

    \item \textbf{Structure of $\G$.} The capability space requires structure beyond a bare set: at minimum a $\sigma$-algebra with a measure so that $\vol(G)$ is well-defined. Richer structure---a lattice with a natural subsumption order, or a topological space where nearby capabilities are similar---would enable stronger results. This choice determines what ``volume'' means and remains the most urgent formal question.

    \item \textbf{Measurement.} What sensors and signals does a GFM actor need to construct the local volume estimator $\Vhat$ in practice? The estimator requires observable capability changes across agents, but the mapping from real-world observations to capability-space signals is unspecified.

    \item \textbf{Convergence.} Does the local finite-difference estimator converge to the true $\vol(G)$ trajectory over time? Under what conditions on observation coverage, trust-model accuracy, and population dynamics does the estimator track the ground truth rather than drifting?

    \item \textbf{Plurality.} Can multiple GFM actors with different models of $\G$ coordinate? When actors disagree about the structure of the capability space itself---not just the shape of $G$ within it---the shared-objective argument breaks down and requires extension.

    \item \textbf{Bootstrap.} Initializing a GFM actor requires estimating other agents' capability sets, which is itself a hard inference problem made harder by the need to discover the structure of $\G$ simultaneously. The goal estimation problem and the capability-space structure problem are coupled, and solving one presupposes progress on the other.
\end{itemize}
